% !TeX root = ../prompt-engineering-guide-for-beginners.tex

%%%%%%%%%%%%%%%%%%%%%%%%%%%%%%%%%%
%% 第十二章：安全与伦理
%%%%%%%%%%%%%%%%%%%%%%%%%%%%%%%%%%
\chapter{安全与伦理：负责任地使用 AI}
\label{ch:safety-ethics}

AI 是一个强大的工具，但强大也意味着需要谨慎使用。
这一章不是要吓唬你，而是帮你了解使用 AI 时的「红线」和「灰色地带」，让你用得安心、用得明白。

\section{隐私保护：哪些内容不该输入 AI}

很多人在使用 AI 时，会不假思索地把各种信息贴进去。比如把一封包含客户联系方式的邮件直接粘贴给 AI 帮忙润色，或者把公司内部的未公开财务报表发给 AI 做分析。这些操作看似方便，但你需要知道一个技术常识：\textbf{在默认设置下，你输入给公共 AI 模型（如免费版 ChatGPT、Claude）的所有 Prompt 和文件，都有可能被用作语料库，参与下一代模型的训练}。一旦被吸收进模型权重（Weights）中，这些机密信息就有可能在别人提问时被“吐”出来。

\textbf{绝对不能直接输入公共 AI 的内容清单（数据红线）：}

\begin{itemize}
  \item \textbf{个人敏感信息}：身份证号、银行卡号、密码、手机号、家庭住址等
  \item \textbf{公司机密}：内部财务数据、未发布的产品计划、客户名单、商业合同等
  \item \textbf{他人隐私}：同事的聊天记录、客户的个人信息、他人的医疗或法律记录
  \item \textbf{受保护的数据}：受保密协议（NDA）约束的内容、政府涉密信息
\end{itemize}

\textbf{企业级数据安全与脱敏建议：}

\begin{itemize}
  \item \textbf{建立本地脱敏机制（Data Masking）}：在将数据发送给云端 AI 之前，先将真实姓名、具体金额、核心项目代号替换为空占位或虚拟标签（如 \verb|[Client_A]|、\verb|[Amount_X]|），等云端推理并返回格式化文本后，再通过本地工具一键映射替换回来。这是最实用的隔离手段。
  \item \textbf{采用具备数据隔离（Zero Data Retention）承诺的服务}：使用正规云服务提供商的 API（如 OpenAI API、Azure OpenAI、AWS Bedrock）或购买认证的企业版（Enterprise/Team 版）。它们在用户协议中明确承诺：API 调用期间既不会缓存用户会话记录，也不会收集这批数据用于基础模型的持续训练。
  \item \textbf{部署本地私有化模型（Local LLMs）}：对于涉密或极高隐私敏感度的工作由于不能连通外网，最佳选择是依靠诸如 Llama 3、Qwen 这样高性能的开源模型，部署在本地电脑或内网机房进行纯离线推理。
  \item \textbf{一个简单的判断标准}：\textbf{绝不将任何你不敢直接发给陌生人或公开在社交平台上的机要信息贴进网页端免费版 AI 的对话框中}。
\end{itemize}

\begin{tipbox}
各大 AI 主流平台（如 ChatGPT 的 Data Controls 设置中）通常都配有“聊天记录与训练（Chat History \& Training）”或类似开关。使用前请务必确认已经关闭了“参与模型提升”服务，尽量从源头消除信息泄漏风险。
\end{tipbox}

\section{事实核查：AI 说的不一定是对的}

前面的章节多次提到了 AI 的「幻觉」问题。这里要再次强调：\textbf{AI 的回答不能作为唯一的事实依据}。这不是否认 AI 的价值，而是帮你在正确的场景中使用它。

以下场景中需要特别警惕：

\begin{itemize}
  \item \textbf{医疗健康}：AI 不是医生。它可以帮你了解症状的一般信息，但不能替代专业诊断。身体不舒服，请去医院。
  \item \textbf{法律咨询}：AI 不了解你所在地区的具体法律法规的最新修订。涉及法律决策时，请咨询专业律师。
  \item \textbf{财务投资}：AI 给的投资建议没有考虑到你的具体财务状况和风险承受能力，也不能预测市场走势。
  \item \textbf{学术引用}：AI 可能会编造论文标题、作者和期刊名。如果你需要引用文献，一定要手动验证每一条引用是否真实存在。
\end{itemize}

\textbf{核心原则（Human-in-the-loop）}：始终坚守“人类控制回路系统”的角色定位。把 AI 当作「不知疲倦的初级研究员」或「初稿的草拟者」，而不是「真理机器」或「全科专家」。它是一个概率推断模型而非精准数据库。只要是在可能带来经济损失、人身威胁、声誉风险的业务决策中，你永远应当作为「最后一环」承担核验与兜底。

\section{版权与原创性}

AI 生成的内容在版权方面目前还有很多不确定性。不同国家和地区的法律对此有不同的看法，而且法律还在不断演变中。这个领域目前还没有明确的全球统一标准，但有几个关键点值得你了解：

\begin{itemize}
  \item \textbf{AI 生成的内容不一定享有版权保护}：在一些法律体系中，只有人类创作的作品才能受到版权保护。纯粹由 AI 生成的内容可能不被视为「作品」。
  \item \textbf{AI 可能无意中复制训练数据}：虽然概率很低，但 AI 有可能生成与训练数据中某段内容非常相似的文字。在发布前最好检查一下。
  \item \textbf{标注 AI 辅助创作是一个好习惯}：如果你的内容有很大比例是由 AI 生成的，标注一下可以体现诚信，也能避免潜在的争议。
\end{itemize}

\textbf{创作权属归属建议}：与其将 AI 视作代笔工具，不如将其设定为“大脑外脑”或“头脑风暴伴侣”。保留对核心思想的原创性把控，让它草拟思路大纲、润色语法结构或寻找论据支撑，而文脉的节奏感和独特价值观点则交由人类自行创作。如果在商业化文本中使用了大范围 AI 原样生成的内容，适度注明“部分灵感由 AI 在监督下生成”既是职业操守的表现，也能从源头上斩断不必要的版权隐患。

\section{AI 偏见与公平性}

AI 模型是从大量的互联网数据中学习的，而互联网上的内容本身就包含各种偏见（性别偏见、种族偏见、文化偏见等）。这意味着 AI 的输出可能也会反映出这些偏见。

\textbf{需要特别注意的场景：}

\begin{itemize}
  \item \textbf{招聘和人事}：用 AI 筛选简历或撰写职位描述时，检查是否有无意中排斥某些群体的措辞。比如 AI 可能会在职位描述中写出「适合年轻人」这样的表述，无意中排除了年龄较大的候选人。
  \item \textbf{内容创作}：AI 生成的文案可能带有刻板印象。比如让它写一个「优秀工程师」的人物敬述，它可能默认写成男性；写「护士」的故事可能默认写成女性。这些隐含的偏见可能在不经意间传递给你的受众。
  \item \textbf{评估和决策}：不要让 AI 单独做出影响他人的重要决策。AI 可以提供分析和建议，但最终的判断应该由人来做
\end{itemize}

\textbf{系统性防偏见（De-biasing）策略}：所有基于大语料训练的算法必然带有互联网历史数据的思维定式（即所谓的统计学分布偏差），对于涉及人文社科和资源分配的任务，在使用 AI 输出之前，务必要从平权与多元维度人工多审视一遍。在系统提示词或要求中，可常备防止偏见的防御性约束：\verb|“在生成用词时，请确保多文化包容性与性别、年龄中立化，绝对不准带入涉及种族、职业、性别或地域色彩等任何刻板偏见，应当保证平权立场与客观分析视角。”|

\section{学术诚信与职场规范}

不同的学校和公司对 AI 使用有不同的规定。有的鼓励使用，有的限制使用，有的明确禁止在特定场景中使用。

\begin{itemize}
  \item \textbf{学生}：在写作业、论文或考试时，先了解你所在学校对 AI 使用的具体政策。很多学校要求你在使用 AI 辅助时必须声明，有些学校完全禁止在特定作业中使用 AI
  \item \textbf{职场人士}：了解公司对 AI 工具使用的规定。有些公司禁止员工向公共 AI 服务输入公司数据，有些公司提供了经过审批的内部 AI 工具
  \item \textbf{自由职业者}：如果你的客户要求原创内容，使用 AI 辅助创作时最好提前沟通，避免交付后产生纠纷
\end{itemize}

\textbf{总结}：AI 是工具，怎么用取决于你。了解规则、尊重边界、保持诚实，你就能在合规的前提下充分发挥 AI 的价值。毕竟，工具本身没有好坏，关键在于使用它的人。
